\documentclass[border=10pt]{standalone}
\usepackage{fontspec}
\setmainfont{Noto Serif}
\newfontfamily\bengalifont{Noto Serif Bengali}[Script=Bengali]
\usepackage[Latin,Bengali]{ucharclasses}
\setTransitionsFor{Bengali}{\begingroup\bengalifont}{\endgroup}
\usepackage{newunicodechar}
\newunicodechar{।}{{\bengalifont\char"0964}}
\newunicodechar{॥}{{\bengalifont\char"0965}}

\usepackage[table]{xcolor}
\usepackage{booktabs,tabularx,array}
\definecolor{bapilrule}{HTML}{1F2937}
\definecolor{bapilstripe}{HTML}{F3F4F6}
\renewcommand{\arraystretch}{1.28}
\setlength{\tabcolsep}{5pt}

\begin{document}
\begin{minipage}{22cm}
{\large\bfseries সারণি ৪৬।}\ {\large দিগন্তভিত্তিক নির্দেশক দীর্ঘমেয়াদি লক্ষ্য}\par\vspace{7pt}
\small
\rowcolors{2}{bapilstripe}{white}
\begin{tabularx}{\linewidth}{@{}>{\hsize=2.000\hsize\raggedright\arraybackslash}X>{\hsize=0.667\hsize\raggedright\arraybackslash}X>{\hsize=0.733\hsize\raggedright\arraybackslash}X>{\hsize=0.733\hsize\raggedright\arraybackslash}X>{\hsize=0.800\hsize\raggedright\arraybackslash}X>{\hsize=1.067\hsize\raggedright\arraybackslash}X@{}}
\arrayrulecolor{bapilrule}\toprule
\rowcolor{white}\textbf{সূচক} & \textbf{২০২৬} & \textbf{২০২৮} & \textbf{২০৩০} & \textbf{২০৩৫} & \textbf{২০৪১} \\
\midrule
পূর্ণ উৎস-তথ্যসহ সূচিবদ্ধ উন্মুক্ত নথি & ০ & ১৫,০০০ & ৩০,০০০ & ১,০০,০০০ & বলবৎ সব উন্মুক্ত দলিল \\
বলবৎ জাতীয় আইনের যে অংশ উদ্ধৃতিসহ বাংলা ও ইংরেজিতে অনুসন্ধানযোগ্য & ০\% & ৫০\% & ৯০\% & ১০০\% & ১০০\%, ৭ দিনের মধ্যে হালনাগাদ \\
BanglaPolicyQA প্রশ্ন (সব সংস্করণ) & ০ & ১,০০০ & ২,৫০০ & ৬,০০০ & বহু-ক্ষেত্রভিত্তিক, রক্ষিত \\
প্রশিক্ষিত কর্মকর্তা (ক্রমপুঞ্জিত) & ০ & ৩০০ & ১,০০০ & ৫,০০০ & জনপ্রশাসন প্রশিক্ষণের অংশ \\
BAPIL-এর সরঞ্জাম বা পদ্ধতি ব্যবহারকারী সংস্থা & ০ & ৩ & ১০ & ৩০ & অধিকাংশ মন্ত্রণালয় ও বিভাগ \\
নীতি অনুসন্ধানের মাসিক সক্রিয় ব্যবহারকারী & ০ & ১০০ & ৫০০ & ৩,০০০ & সরকারজুড়ে নিয়মিত ব্যবহার \\
বাংলা ও ইংরেজির Recall@10 ব্যবধান (পয়েন্ট) & প্রযোজ্য নয় & ৫ বা কম & ৩ বা কম & ২ বা কম & ২ বা কম \\
জনসমক্ষে স্বচ্ছতা নিবন্ধনে তালিকাভুক্ত এআই ব্যবস্থা & ০ & ৬ & ১৫ & সরকারি খাতজুড়ে নিবন্ধন প্রস্তাবিত & প্রমিত চর্চা \\
প্রকাশিত স্বাধীন মূল্যায়ন & ০ & ১ & ৩ & ৮ & অন্তত প্রতি ৫ বছরে \\
DOI-সহ সংরক্ষিত উন্মুক্ত প্রকাশ & ৩ & ৮ & ১৫ & ৪০ & নিরবচ্ছিন্ন \\
\bottomrule
\end{tabularx}
\par\vspace{5pt}{\footnotesize\itshape পরিকল্পনা-লক্ষ্য, প্রতিশ্রুতি নয়। এগুলো সারণি ৪৭-এর ভিত্তি-দৃশ্যকল্প ধরে নেয় এবং প্রতিটি কৌশলগত পর্যালোচনায় পুনর্নির্ধারিত হয়।}
\par\vspace{3pt}{\scriptsize\color{gray} উৎস: BAPIL কারিগরি ব্লুপ্রিন্ট, সংস্করণ ১.০ (২০২৬)। CC BY 4.0।}
\end{minipage}
\end{document}
